\chapter{MCP interactive documentation for all Bijou components}\label{record-294}

\noindent{\footnotesize\color{BookMuted}\path{docs/method/retro/DF_mcp-interactive-documentation-for-all-bijou-components.md}\par}\smallskip

\begin{RecordMeta}
title: MCP interactive documentation for all Bijou components
legend: DF
lane: v5.0.0

\end{RecordMeta}

bijou-mcp should serve as interactive documentation for all of Bijou, not just a rendering service. Add a bijou\_docs tool (or MCP resource endpoints) that returns structured, machine-readable documentation for any component: what it does, when to use it, when NOT to use it, interaction profile behavior (interactive/static/pipe/accessible), related components, and example output. Think of it as the design-system field guide but queryable by LLMs --- an agent could ask ``when should I use a table vs a DAG?'' or ``what component handles step progress?'' and get an authoritative answer with examples. This makes Bijou self-documenting inside any MCP-capable coding tool.

\section{Progress (2026-04-12)}\label{record-294-progress-2026-04-12}

\texttt{b\allowbreak{}i\allowbreak{}j\allowbreak{}o\allowbreak{}u\allowbreak{}\_\allowbreak{}d\allowbreak{}o\allowbreak{}c\allowbreak{}s} now exists and already covers the current \texttt{b\allowbreak{}i\allowbreak{}j\allowbreak{}o\allowbreak{}u\allowbreak{}-\allowbreak{}m\allowbreak{}c\allowbreak{}p} render-tool surface. It also documents a first bounded set of public first-party component families that do not yet have dedicated MCP render tools, including \texttt{m\allowbreak{}a\allowbreak{}r\allowbreak{}k\allowbreak{}d\allowbreak{}o\allowbreak{}w\allowbreak{}n\allowbreak{}(\allowbreak{})}, \texttt{g\allowbreak{}u\allowbreak{}i\allowbreak{}d\allowbreak{}e\allowbreak{}d\allowbreak{}F\allowbreak{}l\allowbreak{}o\allowbreak{}w\allowbreak{}(\allowbreak{})}, \texttt{p\allowbreak{}r\allowbreak{}e\allowbreak{}f\allowbreak{}e\allowbreak{}r\allowbreak{}e\allowbreak{}n\allowbreak{}c\allowbreak{}e\allowbreak{}L\allowbreak{}i\allowbreak{}s\allowbreak{}t\allowbreak{}S\allowbreak{}u\allowbreak{}r\allowbreak{}f\allowbreak{}a\allowbreak{}c\allowbreak{}e\allowbreak{}(\allowbreak{})}, the spinner/timer family, and the current data-viz surfaces such as \texttt{s\allowbreak{}p\allowbreak{}a\allowbreak{}r\allowbreak{}k\allowbreak{}l\allowbreak{}i\allowbreak{}n\allowbreak{}e\allowbreak{}(\allowbreak{})} and \texttt{p\allowbreak{}e\allowbreak{}r\allowbreak{}f\allowbreak{}O\allowbreak{}v\allowbreak{}e\allowbreak{}r\allowbreak{}l\allowbreak{}a\allowbreak{}y\allowbreak{}S\allowbreak{}u\allowbreak{}r\allowbreak{}f\allowbreak{}a\allowbreak{}c\allowbreak{}e\allowbreak{}(\allowbreak{})}.

The remaining gap is breadth, not direction: the catalog still is not a complete field guide for every first-party component or every DOGFOOD surface, and the docs-only families still stop short of a one-tool-per-component interactive render story.

\section{Completed (2026-04-14)}\label{record-294-completed-2026-04-14}

\texttt{b\allowbreak{}i\allowbreak{}j\allowbreak{}o\allowbreak{}u\allowbreak{}\_\allowbreak{}d\allowbreak{}o\allowbreak{}c\allowbreak{}s} now covers the public first-party Bijou component-family surface exposed by \texttt{@\allowbreak{}f\allowbreak{}l\allowbreak{}y\allowbreak{}i\allowbreak{}n\allowbreak{}g\allowbreak{}r\allowbreak{}o\allowbreak{}b\allowbreak{}o\allowbreak{}t\allowbreak{}s\allowbreak{}/\allowbreak{}b\allowbreak{}i\allowbreak{}j\allowbreak{}o\allowbreak{}u}, including the previously missing \texttt{n\allowbreak{}o\allowbreak{}t\allowbreak{}e\allowbreak{}(\allowbreak{})}, the core choice/form families, expressive branding helpers, and the \texttt{r\allowbreak{}e\allowbreak{}n\allowbreak{}d\allowbreak{}e\allowbreak{}r\allowbreak{}B\allowbreak{}y\allowbreak{}M\allowbreak{}o\allowbreak{}d\allowbreak{}e\allowbreak{}(\allowbreak{})} authoring helper. The catalog still prefers real MCP render tools when they exist, and it now synthesizes representative plain-text examples for the remaining docs-only families.

